% cv.tex -- main CV document. Hand-maintained.
%
% Compiles standalone with pdflatex (no biblatex/biber). Publication and
% patent entries are NOT hardcoded here: they come entirely from
% % publications.tex -- GENERATED by cv/scripts/render.py
% Do not hand-edit. Source data: publications.bib + overrides.yml.
% Every \begin{pubcounter} resumes the SAME named counter series
% (defined once in cv.sty), so numbering stays continuous across
% every section below, in the order they appear here.

\subsection*{First-Authored Publications}
\begin{pubcounter}
\item \textbf{Jianxiang Zhao}, Y. He, Y. Wang, L. Nwaeto, X. Chen, T. Cooper, M. Kheradmandkeysomi, Chul B. Park. Bio-based shape memory foams uniting passive daytime radiative cooling with adaptive thermal insulation. \textit{Advanced Functional Materials}, 2026. \textit{(To be submitted)}
\item \textbf{Jianxiang Zhao}, Abdullah Al Faysal, Jun Uk Lee, Tobin Filleter, Patrick C. Lee. Toughening mechanism in micro-layered polymer films revealed by in situ micro-tensile testing. \textit{ACS Applied Materials \& Interfaces}, 2026. \textit{(Under review)}
\item \textbf{Jianxiang Zhao}, Lei Zhang, Jun Uk Lee, Nello D. Sansone, Janet Kim, Leen Bazbaz, Abdullah Al Faysal, Patrick C. Lee. Butterfly Wing Microstructure Inspired Solid/Porous Alternating Layered Structures: In Situ Visualization of Confined Foaming. \textit{Small}, 2025, 21(40). \url{https://doi.org/10.1002/smll.202508472}
\item \textbf{Jianxiang Zhao}, Hongbo Li, Duyoung Choi, Patrick C. Lee. Micro-layered film and foam alternating structures: A novel passive daytime radiative cooling material. \textit{Nano Energy}, 2024, 127, 109695. \url{https://doi.org/10.1016/j.nanoen.2024.109695}
\end{pubcounter}

\subsection*{Co-Authored Publications}
\begin{pubcounter}
\item Isabela Braga, \textbf{Jianxiang Zhao}, P. Mapa, J. Kim, Y. V. G. Ferreira, Rafaela Aguiar, F. G. F. Coelho, Patrick C. Lee, A. P. Braga. An automated image-processing framework for bubble tracking in multilayer polymer foams. \textit{Engineering Applications of Artificial Intelligence}, 2026. \textit{(Under review)}
\item Abdullah Al Faysal, \textbf{Jianxiang Zhao}, Patrick C. Lee. Layered foam/film polymer composite structure exhibiting high absorption-dominant electromagnetic interference (EMI) shielding properties. Americas Regional Meeting of the Polymer Processing Society, Guelph, Canada, 2025.
\item Abdullah Al Faysal, Giulia Mallamaci, \textbf{Jianxiang Zhao}, Cyrille Sollogoub, Patrick C. Lee. Multi-interphase biodegradable nanolayer films for enhanced mechanical and gas barrier properties. \textit{Sustainable Materials and Technologies}, 2026. \textit{(To be submitted)}
\item Taeyoung Lee, Jun Uk Lee, \textbf{Jianxiang Zhao}, Patrick C. Lee. Artificial nacre materials made by CO2 laser-induced technology for graphene structures with stacked PBAT films embedded with cellulose-rich biofibers. \textit{Advanced Functional Materials}, 2026. \textit{(To be submitted)}
\item Haoyu Ma, Maryam Fashandi, Zeineb Ben Rejeb, Piyapong Buahom, \textbf{Jianxiang Zhao}, Pengjian Gong, Qiwu Shi, Guangxian Li, Chul B. Park. POSS/PVTMS aerogels for passive cooling and THz communication
<i>via</i>
cross-linking density regulation and nanoscale bimodal design. \textit{Journal of Materials Chemistry A}, 2024, 12(16), 9627--9636. \url{https://doi.org/10.1039/d4ta00315b}
\item Sungwoong Choi, \textbf{Jianxiang Zhao}, Patrick C. Lee, Duyoung Choi. The Effect of Coupling Agents and Graphene on the Mechanical Properties of Film-Based Post-Consumer Recycled Plastic. \textit{Polymers}, 2024, 16(3), 380. \url{https://doi.org/10.3390/polym16030380}
\item Troy Su, Abdullah Al Faysal, \textbf{Jianxiang Zhao}, Patrick C. Lee. Nano-Layered Films and Foams. In \textit{Polyester Films: Materials, Processes and Applications}, Wiley, 1st ed., pp. 225--276, 2023. Chapter 6. ISBN 9781119535751.
\item Xing Huang, Lei Wang, \textbf{Jianxiang Zhao}, Cong Liu, Fanghao Zhou, Zhongbin Xu, Hao Pei. Fabrication and application of flexible rectangular microtubes. \textit{Polymer-Plastics Technology and Materials}, 2020, 59(9), 959--965.
\end{pubcounter}

\subsection*{Conference Presentations}
\begin{pubcounter}
\item \textbf{Jianxiang Zhao}, Patrick C. Lee. Bio-mimetic micro-/nano-layered foams: In situ visualization of cell nucleation and growth dynamics. Society of Plastics Engineers, ANTEC 2025, Philadelphia, PA, USA, 2025.
\item \textbf{Jianxiang Zhao}, Hongbo Li, Patrick C. Lee. Micro/nano-layered film/foam structures: A novel approach to passive cooling building materials. Society of Plastics Engineers, ANTEC 2024, St. Louis, MO, USA, 2024.
\item \textbf{Jianxiang Zhao}, Hongbo Li, Patrick C. Lee. Micro/nano-layer structured foams with tunable foam morphology and properties. Society of Plastics Engineers, ANTEC 2023, Denver, CO, USA, 2023.
\item \textbf{Jianxiang Zhao}, Patrick C. Lee. High thermal insulation structures using micro-/nano-layered film/foam technology. 37th International Conference of the Polymer Processing Society (PPS-37), Fukuoka, Japan, 2022.
\end{pubcounter}

\subsection*{Patents}
\begin{pubcounter}
\item \textbf{Jianxiang Zhao}. A movable-block variable-differentiation-cross section extrusion die for preparing micro-nano structures. China Patent CN110253856A, 2019.
\item \textbf{Jianxiang Zhao}. A method for generating foamed microcapillary films with a dual network of interconnected pores and internal channels. PCT Patent PCT/CN2019/096518, 2019.
\end{pubcounter}
, which is generated by scripts/render.py from
% publications.bib + overrides.yml. Run `make render` after updating
% manual.bib or overrides.yml, then `make pdf`.
\documentclass[11pt]{article}
\usepackage{cv}

\begin{document}

\cvheader{Jianxiang Zhao}{Curriculum Vitae}%
  {Postdoctoral Fellow \quad Department of Mechanical \& Industrial Engineering \quad University of Toronto}%
  {5 King's College Road, Toronto, ON, Canada}%
  {\href{mailto:jx.zhao@mail.utoronto.ca}{jx.zhao@mail.utoronto.ca} \textbar\ \href{mailto:cv@jxzhao.com}{cv@jxzhao.com} \textbar\ \href{https://jxzhao.com}{jxzhao.com}}%
  {\href{https://orcid.org/0000-0002-3453-2254}{ORCID: 0000-0002-3453-2254} \textbar\ \href{https://scholar.google.com/citations?user=xoSH_fMAAAAJ}{Google Scholar}}

\section{Education}
\textbf{PhD, Mechanical Engineering} \\
University of Toronto, Canada \textbar\ 2021--2025

\textbf{MSc, Chemical Process Equipment} \\
Zhejiang University, China \textbar\ 2017--2020

\textbf{BSc, Control Engineering} \\
Xi'an Jiaotong University, China \textbar\ 2013--2017

\section{Research Interests}
Multifunctional polymeric materials; advanced manufacturing techniques; micro/nano-scale structures; porous structures; physics-constrained machine learning; thermal management and passive cooling.

\section{Technical Expertise}

\subsection{Micro/Nano Structuring \& Fabrication}
Foaming of micro/nano cellular structures; micro/nano-layered film and foam structures; nanofibrillation (fiber-in-fiber) processing; porous microcapillary film extrusion; multilayer coextrusion; gas-assisted microextrusion; CO\textsubscript{2} and UV laser micro/nanostructuring (LIFT).

\subsection{Functional Device \& System Design}
Triboelectric nanogenerator (TENG) fabrication; passive daytime radiative cooling materials; thermal management and insulation structures; EMI shielding composites; shape-memory polymer actuator systems; flexible and stretchable substrate processing.

\subsection{Characterization \& Analysis}
Scanning electron microscopy (operation, sample preparation, and routine instrument maintenance); optical microscopy and in situ visualization of foaming; in situ micro-tensile testing; thermal analysis (DSC, TGA); UV--Vis--NIR and FTIR spectroscopy for solar reflectance and infrared emissivity; thermal conductivity and steady-state insulation measurement; quantitative image-based morphology analysis of cellular structures.

\section{Teaching \& Mentoring}

\subsection{Graduate Student Mentoring}
\begin{itemize}
  \item Maichel Tawdrous, MEng, University of Toronto (2022).
  \item Wesam Rabba, MEng, University of Toronto (2024).
  \item Utkarsh Chadha, MEng, University of Toronto (2024).
\end{itemize}

\subsection{Undergraduate Research Mentoring}
\begin{itemize}
  \item Chengyao Yang, summer research project, University of Toronto (2022).
  \item Chris (Hau) Yung, summer research project, University of Toronto (2023).
  \item Leen Bazbaz, summer research project, University of Toronto (2024). Co-author on \emph{Small} 21.40 (2025): e08472.
\end{itemize}

\subsection{Laboratory and Instrument Training}
Trained graduate and undergraduate students in laboratory safety procedures, covering sample preparation, independent operation, and routine maintenance.

\section{Publications \& Patents}
% Continuous numbering across first-author, co-authored, conference, and
% patent sections is enforced by the shared `pubcounter` list series defined
% in cv.sty; publications.tex resumes that series in every subsection below.
% publications.tex -- GENERATED by cv/scripts/render.py
% Do not hand-edit. Source data: publications.bib + overrides.yml.
% Every \begin{pubcounter} resumes the SAME named counter series
% (defined once in cv.sty), so numbering stays continuous across
% every section below, in the order they appear here.

\subsection*{First-Authored Publications}
\begin{pubcounter}
\item \textbf{Jianxiang Zhao}, Y. He, Y. Wang, L. Nwaeto, X. Chen, T. Cooper, M. Kheradmandkeysomi, Chul B. Park. Bio-based shape memory foams uniting passive daytime radiative cooling with adaptive thermal insulation. \textit{Advanced Functional Materials}, 2026. \textit{(To be submitted)}
\item \textbf{Jianxiang Zhao}, Abdullah Al Faysal, Jun Uk Lee, Tobin Filleter, Patrick C. Lee. Toughening mechanism in micro-layered polymer films revealed by in situ micro-tensile testing. \textit{ACS Applied Materials \& Interfaces}, 2026. \textit{(Under review)}
\item \textbf{Jianxiang Zhao}, Lei Zhang, Jun Uk Lee, Nello D. Sansone, Janet Kim, Leen Bazbaz, Abdullah Al Faysal, Patrick C. Lee. Butterfly Wing Microstructure Inspired Solid/Porous Alternating Layered Structures: In Situ Visualization of Confined Foaming. \textit{Small}, 2025, 21(40). \url{https://doi.org/10.1002/smll.202508472}
\item \textbf{Jianxiang Zhao}, Hongbo Li, Duyoung Choi, Patrick C. Lee. Micro-layered film and foam alternating structures: A novel passive daytime radiative cooling material. \textit{Nano Energy}, 2024, 127, 109695. \url{https://doi.org/10.1016/j.nanoen.2024.109695}
\end{pubcounter}

\subsection*{Co-Authored Publications}
\begin{pubcounter}
\item Isabela Braga, \textbf{Jianxiang Zhao}, P. Mapa, J. Kim, Y. V. G. Ferreira, Rafaela Aguiar, F. G. F. Coelho, Patrick C. Lee, A. P. Braga. An automated image-processing framework for bubble tracking in multilayer polymer foams. \textit{Engineering Applications of Artificial Intelligence}, 2026. \textit{(Under review)}
\item Abdullah Al Faysal, \textbf{Jianxiang Zhao}, Patrick C. Lee. Layered foam/film polymer composite structure exhibiting high absorption-dominant electromagnetic interference (EMI) shielding properties. Americas Regional Meeting of the Polymer Processing Society, Guelph, Canada, 2025.
\item Abdullah Al Faysal, Giulia Mallamaci, \textbf{Jianxiang Zhao}, Cyrille Sollogoub, Patrick C. Lee. Multi-interphase biodegradable nanolayer films for enhanced mechanical and gas barrier properties. \textit{Sustainable Materials and Technologies}, 2026. \textit{(To be submitted)}
\item Taeyoung Lee, Jun Uk Lee, \textbf{Jianxiang Zhao}, Patrick C. Lee. Artificial nacre materials made by CO2 laser-induced technology for graphene structures with stacked PBAT films embedded with cellulose-rich biofibers. \textit{Advanced Functional Materials}, 2026. \textit{(To be submitted)}
\item Haoyu Ma, Maryam Fashandi, Zeineb Ben Rejeb, Piyapong Buahom, \textbf{Jianxiang Zhao}, Pengjian Gong, Qiwu Shi, Guangxian Li, Chul B. Park. POSS/PVTMS aerogels for passive cooling and THz communication
<i>via</i>
cross-linking density regulation and nanoscale bimodal design. \textit{Journal of Materials Chemistry A}, 2024, 12(16), 9627--9636. \url{https://doi.org/10.1039/d4ta00315b}
\item Sungwoong Choi, \textbf{Jianxiang Zhao}, Patrick C. Lee, Duyoung Choi. The Effect of Coupling Agents and Graphene on the Mechanical Properties of Film-Based Post-Consumer Recycled Plastic. \textit{Polymers}, 2024, 16(3), 380. \url{https://doi.org/10.3390/polym16030380}
\item Troy Su, Abdullah Al Faysal, \textbf{Jianxiang Zhao}, Patrick C. Lee. Nano-Layered Films and Foams. In \textit{Polyester Films: Materials, Processes and Applications}, Wiley, 1st ed., pp. 225--276, 2023. Chapter 6. ISBN 9781119535751.
\item Xing Huang, Lei Wang, \textbf{Jianxiang Zhao}, Cong Liu, Fanghao Zhou, Zhongbin Xu, Hao Pei. Fabrication and application of flexible rectangular microtubes. \textit{Polymer-Plastics Technology and Materials}, 2020, 59(9), 959--965.
\end{pubcounter}

\subsection*{Conference Presentations}
\begin{pubcounter}
\item \textbf{Jianxiang Zhao}, Patrick C. Lee. Bio-mimetic micro-/nano-layered foams: In situ visualization of cell nucleation and growth dynamics. Society of Plastics Engineers, ANTEC 2025, Philadelphia, PA, USA, 2025.
\item \textbf{Jianxiang Zhao}, Hongbo Li, Patrick C. Lee. Micro/nano-layered film/foam structures: A novel approach to passive cooling building materials. Society of Plastics Engineers, ANTEC 2024, St. Louis, MO, USA, 2024.
\item \textbf{Jianxiang Zhao}, Hongbo Li, Patrick C. Lee. Micro/nano-layer structured foams with tunable foam morphology and properties. Society of Plastics Engineers, ANTEC 2023, Denver, CO, USA, 2023.
\item \textbf{Jianxiang Zhao}, Patrick C. Lee. High thermal insulation structures using micro-/nano-layered film/foam technology. 37th International Conference of the Polymer Processing Society (PPS-37), Fukuoka, Japan, 2022.
\end{pubcounter}

\subsection*{Patents}
\begin{pubcounter}
\item \textbf{Jianxiang Zhao}. A movable-block variable-differentiation-cross section extrusion die for preparing micro-nano structures. China Patent CN110253856A, 2019.
\item \textbf{Jianxiang Zhao}. A method for generating foamed microcapillary films with a dual network of interconnected pores and internal channels. PCT Patent PCT/CN2019/096518, 2019.
\end{pubcounter}


\section{Professional Service}
Peer reviewer for international journals, 20+ manuscripts reviewed, including \emph{Chemical Engineering Journal}, \emph{Polymer-Plastics Technology and Materials}, and \emph{Progress in Additive Manufacturing}.

\end{document}
